\documentclass[11pt]{article}
\usepackage{kctdocs}
\begin{document}
\kcttitle{05}{Backend Schema}{%
Status & Draft v1 \\
Database & SQLite 3 (WAL mode) via SQLAlchemy 2 + Alembic. Portable to PostgreSQL without schema changes (no SQLite-only types; JSON is stored as TEXT with CHECK \texttt{json\_\allowbreak{}valid}, which becomes JSONB on Postgres) \\
Related & \href{02_TRD.pdf}{02 TRD}, \href{04_APP_FLOW.pdf}{04 App Flow} \\
}
{\small\setlength{\parskip}{0pt}\setcounter{tocdepth}{1}\tableofcontents}
\bigskip
\hypertarget{entity-overview}{%
\section{Entity overview}\label{entity-overview}}

\begin{figure}[!htbp]\centering
\resizebox{\ifdim\width>\linewidth\linewidth\else\width\fi}{!}{\begin{tikzpicture}[x=1cm,y=1cm]
  \pgfdeclarelayer{edges}\pgfsetlayers{edges,main}
  \foreach \n/\x/\y in {
    recordings/5.2/0, test_set_recordings/0.9/0, jobs/13.8/0,
    source_transcripts/0.9/-1.5, alignment_runs/5.2/-1.5, segments/9.5/-1.5, transcriptions/13.8/-1.5,
    transcript_turns/0.9/-3.0, aligned_words/5.2/-3.0, segment_turns/9.5/-3.0, transcription_segments/13.8/-3.0,
    hypotheses/5.2/-4.5, segment_revisions/9.5/-4.5, model_versions/13.8/-4.5,
    language_spans/5.2/-6.0, dataset_items/9.5/-6.0, training_runs/13.8/-6.0,
    eval_results/5.2/-7.5, datasets/9.5/-7.5, eval_runs/13.8/-7.5}
    \node[ktable] (\n) at (\x,\y) {\detokenize\expandafter{\n}};
  \begin{pgfonlayer}{edges}
  \foreach \c/\p in {
    test_set_recordings/recordings, source_transcripts/recordings, alignment_runs/recordings,
    segments/recordings, transcript_turns/source_transcripts, alignment_runs/source_transcripts,
    aligned_words/alignment_runs, aligned_words/transcript_turns, segments/alignment_runs,
    segment_turns/segments, segment_revisions/segments,
    hypotheses/segments, hypotheses/model_versions, language_spans/segment_revisions,
    dataset_items/segment_revisions, dataset_items/datasets, training_runs/datasets,
    model_versions/training_runs, eval_runs/model_versions, eval_runs/datasets,
    transcriptions/jobs, transcriptions/model_versions,
    transcription_segments/transcriptions, segment_revisions/hypotheses}
    \draw[kfk] (\c) -- (\p);
  \draw[kfk] (segment_turns.south) -- ++(0,-0.3) -| (transcript_turns.south);
  \draw[kfk] (eval_results.south) -- ++(0,-0.35) -| (eval_runs.south);
  \draw[kfk] (dataset_items.north east) to[out=60,in=-60] (segments.south east);
  \draw[kfk] (eval_results.north west) to[out=150,in=-150] (segments.west);
  \draw[kfk,dashed] (transcription_segments.west) -- node[klabel] {promoted} (segments.east);
  \end{pgfonlayer}
\end{tikzpicture}
}
\caption{Entity overview. Arrows point from the referencing table to the referenced table.}
\end{figure}

Conventions:

\begin{itemize}
\tightlist
\item
  IDs are ULIDs stored as \texttt{TEXT} (sortable by creation time and safe to generate in either process).
\item
  Timestamps are ISO-8601 UTC strings in \texttt{TEXT} (\texttt{created\_\allowbreak{}at}, \texttt{updated\_\allowbreak{}at}).
\item
  Audio times are \texttt{REAL} seconds relative to the recording start.
\item
  Enumerations are \texttt{TEXT} columns with \texttt{CHECK} constraints; in SQLAlchemy they map to \texttt{Enum(.\allowbreak{}.\allowbreak{}.\allowbreak{},\allowbreak{} native\_\allowbreak{}enum=\allowbreak{}False)}.
\end{itemize}

\hypertarget{ddl}{%
\section{DDL}\label{ddl}}

\begin{lstlisting}[language=SQL]
PRAGMA foreign_keys = ON;
PRAGMA journal_mode = WAL;

-- ───────────────────────── Sources ─────────────────────────
CREATE TABLE recordings (
  id              TEXT PRIMARY KEY,
  title           TEXT NOT NULL,
  original_path   TEXT NOT NULL,            -- data/store/originals/<sha256>.<ext>
  normalized_path TEXT,                     -- data/store/audio16k/<sha256>.wav
  sha256          TEXT NOT NULL UNIQUE,
  duration_s      REAL NOT NULL,
  domain          TEXT NOT NULL DEFAULT 'interview'
                  CHECK (domain IN ('interview','tv','radio','church','meeting','other')),
  consent_scope   TEXT NOT NULL
                  CHECK (consent_scope IN ('training','eval_only','transcription_only')),
  recorded_on     TEXT,
  region          TEXT,                     -- free text, optional; never ethnicity
  status          TEXT NOT NULL DEFAULT 'imported'
                  CHECK (status IN ('imported','aligning','ready_for_review',
                                    'alignment_suspect','reviewed','archived')),
  notes           TEXT,
  created_at      TEXT NOT NULL,
  updated_at      TEXT NOT NULL
);

CREATE TABLE source_transcripts (
  id              TEXT PRIMARY KEY,
  recording_id    TEXT NOT NULL REFERENCES recordings(id) ON DELETE CASCADE,
  path            TEXT NOT NULL,
  format          TEXT NOT NULL CHECK (format IN ('docx','pdf','txt')),
  sha256          TEXT NOT NULL,
  parser          TEXT NOT NULL,            -- 'python-docx' | 'pypdf'
  parser_version  TEXT NOT NULL,
  italic_means    TEXT CHECK (italic_means IN ('sw','en','none')) DEFAULT 'sw',
  created_at      TEXT NOT NULL,
  UNIQUE (recording_id, sha256)
);

CREATE TABLE transcript_turns (
  id              TEXT PRIMARY KEY,
  source_transcript_id TEXT NOT NULL REFERENCES source_transcripts(id) ON DELETE CASCADE,
  ordinal         INTEGER NOT NULL,
  speaker_label   TEXT,                     -- 'I', 'R', 'P1' as written
  speaker_role    TEXT CHECK (speaker_role IN ('interviewer','respondent','host',
                                               'preacher','audience','other')),
  text_raw        TEXT NOT NULL,            -- exactly as in the source (label source)
  text_align      TEXT NOT NULL,            -- normalised text used by the aligner
  word_map        TEXT NOT NULL CHECK (json_valid(word_map)),
                  -- [[align_word_idx, raw_char_start, raw_char_end], ...]
  spans           TEXT NOT NULL DEFAULT '[]' CHECK (json_valid(spans)),
                  -- language hints from formatting: [{start,end,lang}]
  annotations     TEXT NOT NULL DEFAULT '[]' CHECK (json_valid(annotations)),
                  -- [{"marker":"[unclear]","char":123}]
  UNIQUE (source_transcript_id, ordinal)
);

-- ───────────────────────── Alignment ─────────────────────────
CREATE TABLE alignment_runs (
  id              TEXT PRIMARY KEY,
  recording_id    TEXT NOT NULL REFERENCES recordings(id) ON DELETE CASCADE,
  source_transcript_id TEXT NOT NULL REFERENCES source_transcripts(id),
  aligner         TEXT NOT NULL,            -- 'ctc-forced-aligner'
  aligner_version TEXT NOT NULL,
  params          TEXT NOT NULL CHECK (json_valid(params)),
  mean_confidence REAL,
  status          TEXT NOT NULL CHECK (status IN ('running','done','failed')),
  created_at      TEXT NOT NULL,
  finished_at     TEXT
);

CREATE TABLE aligned_words (
  alignment_run_id TEXT NOT NULL REFERENCES alignment_runs(id) ON DELETE CASCADE,
  turn_id         TEXT NOT NULL REFERENCES transcript_turns(id),
  word_idx        INTEGER NOT NULL,         -- index within turn's text_align
  word            TEXT NOT NULL,
  start_s         REAL NOT NULL,
  end_s           REAL NOT NULL,
  score           REAL NOT NULL,            -- 0..1
  PRIMARY KEY (alignment_run_id, turn_id, word_idx)
);

-- ───────────────────────── Segments & review ─────────────────────────
CREATE TABLE segments (
  id              TEXT PRIMARY KEY,
  recording_id    TEXT NOT NULL REFERENCES recordings(id) ON DELETE CASCADE,
  alignment_run_id TEXT REFERENCES alignment_runs(id),
  origin          TEXT NOT NULL CHECK (origin IN ('alignment','vad','inference_correction')),
  ordinal         INTEGER NOT NULL,
  start_s         REAL NOT NULL,            -- mirrors current revision
  end_s           REAL NOT NULL,
  speaker_label   TEXT,
  status          TEXT NOT NULL DEFAULT 'candidate'
                  CHECK (status IN ('candidate','in_review','draft','approved',
                                    'approved_eval_only','rejected')),
  current_revision_id TEXT,                 -- FK added below (circular)
  alignment_conf  REAL,
  auto_flags      TEXT NOT NULL DEFAULT '[]' CHECK (json_valid(auto_flags)),
                  -- ['low_alignment','unclear_marker','overlap_marker',...]
  priority        REAL NOT NULL DEFAULT 0,
  audio_cache_path TEXT,                    -- rendered slice; invalidated on bound change
  created_at      TEXT NOT NULL,
  updated_at      TEXT NOT NULL,
  CHECK (end_s > start_s),
  UNIQUE (recording_id, ordinal)
);
CREATE INDEX ix_segments_queue ON segments (recording_id, status, priority DESC);

CREATE TABLE segment_turns (
  segment_id      TEXT NOT NULL REFERENCES segments(id) ON DELETE CASCADE,
  turn_id         TEXT NOT NULL REFERENCES transcript_turns(id),
  PRIMARY KEY (segment_id, turn_id)
);

CREATE TABLE segment_revisions (          -- append-only
  id              TEXT PRIMARY KEY,
  segment_id      TEXT NOT NULL REFERENCES segments(id) ON DELETE CASCADE,
  revision_no     INTEGER NOT NULL,
  parent_revision_id TEXT REFERENCES segment_revisions(id),
  text            TEXT NOT NULL,
  start_s         REAL NOT NULL,
  end_s           REAL NOT NULL,
  speaker_label   TEXT,
  flags           TEXT NOT NULL DEFAULT '{}' CHECK (json_valid(flags)),
                  -- {"unclear":false,"overlap":false,"noisy":false,
                  --  "music":false,"non_verbatim":false}
  status          TEXT NOT NULL,            -- status set by this revision
  source          TEXT NOT NULL CHECK (source IN ('aligned_transcript','vad_empty',
                                                  'model_hypothesis','human')),
  hypothesis_id   TEXT REFERENCES hypotheses(id),  -- when source=model_hypothesis
  author          TEXT NOT NULL,            -- 'aligner', 'system', or reviewer handle
  created_at      TEXT NOT NULL,
  UNIQUE (segment_id, revision_no)
);

CREATE TABLE language_spans (
  revision_id     TEXT NOT NULL REFERENCES segment_revisions(id) ON DELETE CASCADE,
  char_start      INTEGER NOT NULL,
  char_end        INTEGER NOT NULL,
  lang            TEXT NOT NULL CHECK (lang IN ('en','sw','mixed','luo','other','unclear')),
  PRIMARY KEY (revision_id, char_start),
  CHECK (char_end > char_start)
);

CREATE TABLE hypotheses (
  id              TEXT PRIMARY KEY,
  segment_id      TEXT NOT NULL REFERENCES segments(id) ON DELETE CASCADE,
  model_version_id TEXT NOT NULL REFERENCES model_versions(id),
  text            TEXT NOT NULL,
  words           TEXT CHECK (json_valid(words)),   -- [{w,start,end,p}]
  avg_logprob     REAL,
  no_speech_prob  REAL,
  wer_vs_current  REAL,                     -- vs revision current at creation time
  created_at      TEXT NOT NULL
);

-- ───────────────────────── Datasets, training, models ─────────────────────────
CREATE TABLE datasets (
  id              TEXT PRIMARY KEY,
  name            TEXT NOT NULL UNIQUE,     -- 'kct-ds-v3'
  filter_spec     TEXT NOT NULL CHECK (json_valid(filter_spec)),
  split_spec      TEXT NOT NULL CHECK (json_valid(split_spec)),
  test_set_name   TEXT NOT NULL,            -- 'test_v1' (frozen definition)
  normalizer_version TEXT NOT NULL,
  sha256          TEXT,                     -- hash of sorted (segment_id, revision_id, split)
  frozen          INTEGER NOT NULL DEFAULT 0 CHECK (frozen IN (0,1)),
  export_path     TEXT,
  stats           TEXT CHECK (json_valid(stats)),   -- hours/segments/recordings per split
  created_at      TEXT NOT NULL
);

CREATE TABLE dataset_items (
  dataset_id      TEXT NOT NULL REFERENCES datasets(id) ON DELETE CASCADE,
  segment_id      TEXT NOT NULL REFERENCES segments(id),
  revision_id     TEXT NOT NULL REFERENCES segment_revisions(id),  -- pinned text
  split           TEXT NOT NULL CHECK (split IN ('train','dev','test')),
  PRIMARY KEY (dataset_id, segment_id)
);
CREATE INDEX ix_dataset_items_split ON dataset_items (dataset_id, split);

CREATE TABLE test_set_recordings (         -- frozen test membership, by recording
  test_set_name   TEXT NOT NULL,
  recording_id    TEXT NOT NULL REFERENCES recordings(id),
  PRIMARY KEY (test_set_name, recording_id)
);

CREATE TABLE training_runs (
  id              TEXT PRIMARY KEY,
  dataset_id      TEXT NOT NULL REFERENCES datasets(id),
  base_model      TEXT NOT NULL,            -- 'openai/whisper-small'
  platform        TEXT NOT NULL CHECK (platform IN ('kaggle','local','other')),
  config          TEXT NOT NULL CHECK (json_valid(config)),
  code_commit     TEXT,                     -- git sha of transcribe package used
  notebook_ref    TEXT,                     -- kaggle notebook slug + version
  dev_metrics     TEXT CHECK (json_valid(dev_metrics)),
  started_at      TEXT,
  finished_at     TEXT,
  status          TEXT NOT NULL CHECK (status IN ('registered','failed'))
);

CREATE TABLE model_versions (
  id              TEXT PRIMARY KEY,
  name            TEXT NOT NULL UNIQUE,     -- 'whisper-small-kct-v1'
  family          TEXT NOT NULL CHECK (family IN ('whisper','w2vbert_ctc','legacy_ctc')),
  base_model      TEXT NOT NULL,
  training_run_id TEXT REFERENCES training_runs(id),  -- NULL for zero-shot baselines
  hf_path         TEXT NOT NULL,            -- data/store/models/<name>/hf
  ct2_path        TEXT,                     -- data/store/models/<name>/ct2-int8
  decode_params   TEXT NOT NULL CHECK (json_valid(decode_params)),
  status          TEXT NOT NULL DEFAULT 'candidate'
                  CHECK (status IN ('converting','candidate','production','retired','broken')),
  promoted_at     TEXT,
  promotion_note  TEXT,
  created_at      TEXT NOT NULL
);
CREATE UNIQUE INDEX ux_one_production ON model_versions (status) WHERE status = 'production';

CREATE TABLE eval_runs (
  id              TEXT PRIMARY KEY,
  model_version_id TEXT NOT NULL REFERENCES model_versions(id),
  dataset_id      TEXT NOT NULL REFERENCES datasets(id),
  split           TEXT NOT NULL CHECK (split IN ('dev','test','external')),
  normalizer_version TEXT NOT NULL,
  metrics         TEXT CHECK (json_valid(metrics)),
                  -- {"wer":..,"cer":..,"wer_en":..,"wer_sw":..,"wer_mixed":..,
                  --  "switch_point_er":..,"halluc_per_min":..,"rtf":..,
                  --  "by_domain":{..},"by_flag":{..}}
  status          TEXT NOT NULL CHECK (status IN ('queued','running','done','failed')),
  created_at      TEXT NOT NULL
);

CREATE TABLE eval_results (
  eval_run_id     TEXT NOT NULL REFERENCES eval_runs(id) ON DELETE CASCADE,
  segment_id      TEXT NOT NULL REFERENCES segments(id),
  reference       TEXT NOT NULL,
  hypothesis      TEXT NOT NULL,
  wer             REAL NOT NULL,
  cer             REAL NOT NULL,
  n_ref_words     INTEGER NOT NULL,
  n_switch_points INTEGER NOT NULL,
  n_switch_errors INTEGER NOT NULL,
  PRIMARY KEY (eval_run_id, segment_id)
);

-- ───────────────────────── Jobs & transcription ─────────────────────────
CREATE TABLE jobs (
  id              TEXT PRIMARY KEY,
  kind            TEXT NOT NULL CHECK (kind IN ('parse_transcript','align','vad_segment',
                        'hypothesize','build_dataset','export_dataset','register_model',
                        'convert_model','evaluate','transcribe','export_transcription')),
  payload         TEXT NOT NULL CHECK (json_valid(payload)),
  status          TEXT NOT NULL DEFAULT 'queued'
                  CHECK (status IN ('queued','running','done','failed','cancelled')),
  progress        REAL NOT NULL DEFAULT 0,  -- 0..1
  attempts        INTEGER NOT NULL DEFAULT 0,
  worker_id       TEXT,
  heartbeat_at    TEXT,
  error           TEXT,
  result          TEXT CHECK (json_valid(result)),
  created_at      TEXT NOT NULL,
  started_at      TEXT,
  finished_at     TEXT
);
CREATE INDEX ix_jobs_queue ON jobs (status, created_at);

CREATE TABLE transcriptions (
  id              TEXT PRIMARY KEY,
  job_id          TEXT NOT NULL REFERENCES jobs(id),
  title           TEXT NOT NULL,
  input_path      TEXT NOT NULL,
  sha256          TEXT NOT NULL,
  duration_s      REAL,
  domain          TEXT NOT NULL DEFAULT 'other',
  consent_to_train INTEGER NOT NULL DEFAULT 0 CHECK (consent_to_train IN (0,1)),
  model_version_id TEXT NOT NULL REFERENCES model_versions(id),
  rtf             REAL,
  expires_at      TEXT,                     -- retention; NULL once promoted
  created_at      TEXT NOT NULL
);

CREATE TABLE transcription_segments (
  id              TEXT PRIMARY KEY,
  transcription_id TEXT NOT NULL REFERENCES transcriptions(id) ON DELETE CASCADE,
  ordinal         INTEGER NOT NULL,
  start_s         REAL NOT NULL,
  end_s           REAL NOT NULL,
  text            TEXT NOT NULL,            -- raw model output (immutable)
  words           TEXT CHECK (json_valid(words)),
  avg_logprob     REAL,
  no_speech_prob  REAL,
  edited_text     TEXT,                     -- user edit; exports prefer this
  edited_at       TEXT,
  promoted_segment_id TEXT REFERENCES segments(id),
  UNIQUE (transcription_id, ordinal)
);
\end{lstlisting}

The circular foreign key \texttt{segments.\allowbreak{}current\_\allowbreak{}revision\_\allowbreak{}id → segment\_\allowbreak{}revisions.\allowbreak{}id} is created by Alembic with \texttt{use\_\allowbreak{}alter=\allowbreak{}True}. SQLite does not enforce ALTER-added foreign keys, so the invariant is checked in the service layer and by a test.

\hypertarget{key-invariants-enforced-in-the-service-layer-and-tested}{%
\section{Key invariants (enforced in the service layer and tested)}\label{key-invariants-enforced-in-the-service-layer-and-tested}}

\begin{enumerate}
\def\labelenumi{\arabic{enumi}.}
\tightlist
\item
  \texttt{segment\_\allowbreak{}revisions} is append-only: no UPDATE or DELETE except by cascade when a recording is deleted.
\item
  \texttt{segments.\allowbreak{}\{start\_\allowbreak{}s,\allowbreak{}end\_\allowbreak{}s,\allowbreak{}status,\allowbreak{}speaker\_\allowbreak{}label\}} always equal those of \texttt{current\_\allowbreak{}revision\_\allowbreak{}id}.
\item
  A \texttt{hypotheses} row never changes \texttt{segment\_\allowbreak{}revisions}. Only an explicit reviewer action creates a revision with \texttt{source=\allowbreak{}'model\_\allowbreak{}hypothesis'}.
\item
  \texttt{datasets.\allowbreak{}frozen =\allowbreak{} 1} means the dataset has no inserts, updates or deletes on \texttt{dataset\_\allowbreak{}items}, and \texttt{sha256} is set.
\item
  The train/dev splits contain only recordings with \texttt{consent\_\allowbreak{}scope=\allowbreak{}'training'}. The test split may contain \texttt{training} or \texttt{eval\_\allowbreak{}only}.
\item
  No \texttt{recording\_\allowbreak{}id} appears in more than one split in a dataset, and no recording in \texttt{test\_\allowbreak{}set\_\allowbreak{}recordings} appears in train or dev of any dataset.
\item
  Items with status \texttt{approved\_\allowbreak{}eval\_\allowbreak{}only}, or with flags \texttt{unclear}, \texttt{non\_\allowbreak{}verbatim} or \texttt{overlap}, never go to train (they may go to test, tagged by flag).
\item
  At most one \texttt{model\_\allowbreak{}versions} row has \texttt{status=\allowbreak{}'production'} (partial unique index).
\end{enumerate}

\hypertarget{file-store-layout}{%
\section{File store layout}\label{file-store-layout}}

\begin{lstlisting}
data/
  transcribe.db
  store/
    originals/<sha256>.<ext>          # byte-identical uploads (never modified)
    transcripts/<sha256>.<docx|pdf>
    audio16k/<sha256>.wav             # normalised full recordings
    slices/<segment_id>_<rev>.wav     # cache; safe to delete
    models/<name>/{hf/, ct2-int8/, run_manifest.json}
    uploads/<transcription_id>.<ext>  # retention-managed
  exports/
    <dataset_name>/
      dataset_card.md
      metadata.jsonl                  # see §5
      train/<segment_id>.wav
      dev/<segment_id>.wav
      test/<segment_id>.wav
\end{lstlisting}

The existing \texttt{data/\allowbreak{}projects/\allowbreak{}<name>/\allowbreak{}\{transcript.\allowbreak{}json,\allowbreak{}audio.\allowbreak{}json,\allowbreak{}chunks/\allowbreak{}\}} manifests are migrated once by \texttt{python -\allowbreak{}m transcribe.\allowbreak{}db.\allowbreak{}migrate\_\allowbreak{}json} (see the Implementation Plan, Phase 0).

\hypertarget{export-format-kaggle--hugging-face}{%
\section{Export format (Kaggle / Hugging Face)}\label{export-format-kaggle--hugging-face}}

\texttt{metadata.\allowbreak{}jsonl} has one line per item and follows the \texttt{audiofolder} convention, so the export loads with \texttt{load\_\allowbreak{}dataset("audiofolder",\allowbreak{} data\_\allowbreak{}dir=\allowbreak{}.\allowbreak{}.\allowbreak{}.\allowbreak{})}:

\begin{lstlisting}
{"file_name": "train/01J8Z...wav", "split": "train", "text": "Sasa mimi nilikuwa nimeapply hiyo job lakini they never called back",
 "segment_id": "01J8Z...", "revision_id": "01J90...", "recording_group": "r_7f3a",
 "domain": "interview", "duration": 11.55, "speaker_role": "respondent",
 "flags": [], "lang_spans": [[0, 26, "sw"], [27, 36, "mixed"], [37, 40, "sw"], [41, 70, "en"]]}
\end{lstlisting}

\begin{itemize}
\tightlist
\item
  \texttt{recording\_\allowbreak{}group} is a salted hash of \texttt{recording\_\allowbreak{}id}, so no titles or file names leave the machine.
\item
  \texttt{dataset\_\allowbreak{}card.\allowbreak{}md} records the dataset name, sha256, normaliser version, per-split hours, filter spec, and the consent statement.
\end{itemize}

The Kaggle notebook writes \texttt{run\_\allowbreak{}manifest.\allowbreak{}json}:

\begin{lstlisting}
{"dataset_name": "kct-ds-v3", "dataset_sha256": "…", "base_model": "openai/whisper-small",
 "code_commit": "…", "config": {…}, "best_step": 3500, "dev_wer": 0.231,
 "started_at": "…", "finished_at": "…", "kaggle_notebook": "user/finetune-whisper@v7"}
\end{lstlisting}

\hypertarget{pydantic-api-schemas-core}{%
\section{Pydantic API schemas (core)}\label{pydantic-api-schemas-core}}

\begin{lstlisting}[language=Python]
class RevisionIn(BaseModel):
    base_revision_id: str
    text: str
    start_s: float
    end_s: float
    speaker_label: str | None = None
    flags: dict[Literal["unclear","overlap","noisy","music","non_verbatim"], bool] = {}
    spans: list[tuple[int, int, Literal["en","sw","mixed","luo","other","unclear"]]] = []
    status: Literal["draft","approved","approved_eval_only","rejected"]

class SegmentOut(BaseModel):
    id: str; recording_id: str; ordinal: int
    start_s: float; end_s: float; status: str; priority: float
    auto_flags: list[str]; alignment_conf: float | None
    current_revision: RevisionOut
    hypotheses: list[HypothesisOut]
    audio_url: str

class DatasetIn(BaseModel):
    name: str
    filter_spec: FilterSpec     # recordings, domains, statuses, exclude_flags, roles, min/max duration
    split_spec: SplitSpec       # dev_fraction, seed, group_by: Literal["recording","speaker_cluster"]
    test_set_name: str = "test_v1"
\end{lstlisting}

\hypertarget{migration-from-the-current-json-manifests}{%
\section{Migration from the current JSON manifests}\label{migration-from-the-current-json-manifests}}

\begin{longtable}[]{@{}>{\raggedright\arraybackslash}p{(\linewidth - 4\tabcolsep) * \real{0.371}}>{\raggedright\arraybackslash}p{(\linewidth - 4\tabcolsep) * \real{0.629}}@{}}
\toprule\noalign{}
Current (\texttt{audio.\allowbreak{}json} / \texttt{transcript.\allowbreak{}json}) & New \\
\midrule\noalign{}
\endhead
\bottomrule\noalign{}
\endlastfoot
\texttt{recording\_\allowbreak{}id}, \texttt{source\_\allowbreak{}audio} & \texttt{recordings} (sha256 computed, original copied to store) \\
\texttt{transcript\_\allowbreak{}turns[]} & \texttt{source\_\allowbreak{}transcripts} + \texttt{transcript\_\allowbreak{}turns} (\texttt{text} → \texttt{text\_\allowbreak{}raw}, spans kept) \\
\texttt{segments[]} with \texttt{status=\allowbreak{}candidate} and empty transcript & Dropped. They are VAD candidates that were never reviewed, and re-running alignment gives better segments \\
\texttt{segments[]} with \texttt{status=\allowbreak{}approved} & \texttt{segments} (origin=\texttt{vad}) + revision r1 (\texttt{source=\allowbreak{}human}, status approved) \\
\texttt{asr\_\allowbreak{}hypothesis}, \texttt{wer} & Ignored (no model existed) \\
\end{longtable}
\end{document}
